% RQ5 Introduction: Inter-DAO Cooperation

\subsection{Motivation}

DAOs do not exist in isolation. They operate within ecosystems of interdependent protocols, share user bases, and face common challenges. Public goods that benefit the entire ecosystem---infrastructure, standards, research---are under-provided when each DAO acts independently.

Inter-DAO cooperation has emerged organically:
\begin{itemize}
    \item \textbf{Protocol collaborations}: Joint liquidity provision, shared security
    \item \textbf{Grants coordination}: DAOs co-funding public goods projects
    \item \textbf{Standard setting}: Multi-DAO governance over shared standards
    \item \textbf{Crisis response}: Coordinated action during security incidents
\end{itemize}

Yet formal mechanisms for inter-DAO cooperation remain primitive. Most cross-DAO interactions rely on informal relationships, multi-sig arrangements, or bilateral negotiations. The governance of multi-DAO initiatives is often ad hoc.

Understanding when and how DAOs can effectively cooperate---and what mechanisms enable coordination---is essential as the ecosystem moves toward greater interdependence.

\subsection{Research Question}

This paper investigates:

\begin{quote}
\textbf{RQ5:} How does inter-DAO cooperation affect governance outcomes and resource efficiency?
\end{quote}

Specifically, we examine:
\begin{enumerate}
    \item Under what conditions do DAOs choose to cooperate on joint initiatives?
    \item How do different cooperation structures (symmetric, asymmetric, hub-spoke) affect outcomes?
    \item What mechanisms enable stable cooperation despite divergent preferences?
    \item How does cooperation affect individual DAO governance metrics (turnout, pass rate)?
\end{enumerate}

\subsection{Cooperation Mechanisms}

We study several inter-DAO cooperation mechanisms:

\subsubsection{Joint Proposals}

Proposals requiring approval from multiple DAOs:
\begin{equation}
\text{pass}(p) = \bigwedge_{d \in \text{DAOs}} \text{approve}_d(p)
\end{equation}

\subsubsection{Resource Pooling}

Shared treasury contributions for joint initiatives:
\begin{equation}
\text{pool} = \sum_{d \in \text{DAOs}} \text{contribution}_d
\end{equation}

\subsubsection{Cross-DAO Voting}

Token holders with positions in multiple DAOs voting across boundaries, creating preference alignment.

\subsubsection{Delegation Networks}

Delegates operating across multiple DAOs, transferring information and norms.

\subsection{Contributions}

This paper contributes:

\begin{enumerate}
    \item \textbf{Cooperation taxonomy}: Classification of inter-DAO relationship types
    \item \textbf{Success conditions}: Identification of factors enabling stable cooperation
    \item \textbf{Mechanism design}: Guidelines for joint proposal structures
    \item \textbf{Simulation framework}: Multi-DAO modeling capabilities for future research
\end{enumerate}

\subsection{Paper Organization}

Section~\ref{sec:background} reviews inter-DAO cooperation in practice. Section~\ref{sec:theory} develops formal models of multi-DAO dynamics. Section~\ref{sec:architecture} describes simulation implementation. Section~\ref{sec:methodology} details experimental design. Section~\ref{sec:results} presents findings. Section~\ref{sec:discussion} interprets results and derives practical guidance. Section~\ref{sec:conclusion} concludes.
